\documentclass[11pt,a4paper]{article}
\usepackage[utf8]{inputenc}
\usepackage{amsmath,amssymb,amsthm,mathtools}
\usepackage[margin=2.3cm]{geometry}
\usepackage{enumitem}
\usepackage{booktabs}
\usepackage{xcolor}
\usepackage{manfnt}
\usepackage{hyperref}
\usepackage{microtype}

\hypersetup{colorlinks=true,linkcolor=black,urlcolor=black,citecolor=black}

\newtheorem{theorem}{Theorem}[section]
\newtheorem{lemma}[theorem]{Lemma}
\newtheorem{corollary}[theorem]{Corollary}
\theoremstyle{definition}
\newtheorem{definition}[theorem]{Definition}
\newtheorem{example}[theorem]{Example}

\newcommand{\Z}{\mathbb{Z}}
\newcommand{\N}{\mathbb{N}}
\newcommand{\md}[1]{\ (\mathrm{mod}\ #1)}
\DeclareMathOperator{\lcm}{lcm}

% Book-style notes with a small-caps heading and no boxes or colors.
\newtheoremstyle{note}{\topsep}{\topsep}{\normalfont}{0pt}{\scshape}{.}{0.6em}{\thmname{#1}}
\newtheoremstyle{noteit}{\topsep}{\topsep}{\itshape}{0pt}{\scshape}{.}{0.6em}{\thmname{#1}}
\theoremstyle{noteit}
\newtheorem*{question}{Question}
\theoremstyle{note}
\newtheorem*{intuition}{Intuition}
\newtheorem*{sketch}{Proof sketch}
\newtheorem*{counter}{Canonical counterexample}
\newtheorem*{thumbrule}{Rule of thumb}
\newtheorem*{mistakeX}{Common mistake}
\newenvironment{mistake}{\begin{mistakeX}\leavevmode\marginpar{\raggedleft\vspace*{-0.9\baselineskip}\dbend}\ignorespaces}{\end{mistakeX}}
\setlength{\marginparwidth}{1.6cm}
\reversemarginpar

\setlist{itemsep=2pt,topsep=4pt}
\setlength{\parindent}{0pt}
\setlength{\parskip}{5pt}

\title{\textbf{Number Theory}\\[2pt]\large A detailed guide to divisibility, congruences, Euler and RSA}
\author{Carla González Mina}
\date{Companion to the \emph{Number Theory} videos in \texttt{manim-clases} \,$\cdot$\, main reference Rosen, ch.~4}

\begin{document}
\maketitle

\begin{quotation}\small
\noindent\textsc{How to use this guide.}\enspace
Each topic opens with the \textbf{question} it answers and the \textbf{intuition} behind it, which is what the videos show. Then comes the \textbf{formal statement}, a full \textbf{proof} and a \textbf{worked example}. Most topics close with three short reminders. The \emph{proof sketch} is the argument in a few steps, the \emph{canonical counterexample} shows what breaks when you drop a hypothesis, and the \emph{rule of thumb} tells you when to use the result. Every module ends with exercises and their solutions.

All the numbers in this guide were checked in Python with \texttt{pow}, \texttt{math.gcd} and exhaustive search.
\end{quotation}

\tableofcontents
\newpage

%======================================================================
\section{Module 1. Divisibility, the division algorithm, gcd and the FTA}
%======================================================================

\subsection{Divisibility}

\begin{question}
What does it mean for one number to fit into another a whole number of times, and which rules can you use without a second thought?
\end{question}

\begin{definition}
Let $a,b\in\Z$ with $a\neq 0$. We say that $a$ \textbf{divides} $b$, written $a\mid b$, if $b=a\cdot k$ for some $k\in\Z$. If there is no such $k$ we write $a\nmid b$.
\end{definition}

\begin{intuition}
$a\mid b$ means you can cover $b$ exactly with blocks of size $|a|$ and have nothing left over. So $3\mid 12$ because $12$ is four blocks of $3$, while $5\nmid 12$ because blocks of 5 leave 2 behind. Almost everything in this guide grows out of one habit, which is to \emph{write what you know as $b=ak$ and then do algebra}.
\end{intuition}

\begin{theorem}[Basic properties]
For all $a,b,c,x,y\in\Z$, with nonzero divisors,
\begin{enumerate}
  \item $1\mid a$, $a\mid 0$, and $a\mid a$ when $a\ne 0$.
  \item \textbf{Transitivity.} If $a\mid b$ and $b\mid c$ then $a\mid c$.
  \item \textbf{Linear combinations.} If $a\mid b$ and $a\mid c$ then $a\mid (bx+cy)$.
  \item \textbf{Size.} If $b\ne 0$ and $a\mid b$ then $|a|\le|b|$.
  \item \textbf{Antisymmetry.} If $a,b>0$, $a\mid b$ and $b\mid a$, then $a=b$. In $\Z$ the conclusion is $a=\pm b$.
\end{enumerate}
\end{theorem}

\begin{proof}
\textbf{(1)} $a=1\cdot a$, $0=a\cdot 0$ and $a=a\cdot 1$.

\textbf{(2)} Write $b=ak$ and $c=b\ell$ with $k,\ell\in\Z$. Then $c=(ak)\ell=a(k\ell)$ and $k\ell$ is an integer.

\textbf{(3)} Write $b=ak$ and $c=a\ell$. Then $bx+cy=akx+a\ell y=a(kx+\ell y)$, and $kx+\ell y$ is an integer.

\textbf{(4)} Write $b=ak$. Here $k\ne 0$, since otherwise $b=0$. A nonzero integer has $|k|\ge 1$, so $|b|=|a||k|\ge|a|$.

\textbf{(5)} By (4), $a\le b$ and $b\le a$, so $a=b$. Over $\Z$ the same argument with absolute values gives $|a|=|b|$.
\end{proof}

Property (3) is the tool you'll use most. A huge number of proofs end with ``$d$ divides both terms, so it divides their sum''.

\begin{counter}
The converse of (3) is false. $3\mid (1+2)$, yet $3\nmid 1$ and $3\nmid 2$. Knowing that $d$ divides a sum tells you nothing about the individual terms.
\end{counter}

\begin{thumbrule}
To prove $a\mid(\text{something})$, write that something as $a$ times an integer, or as a combination of things $a$ already divides.
\end{thumbrule}

\subsection{The division algorithm}

\begin{question}
When $d\nmid a$, how do we describe exactly what's left over, and why is there only one way to do it?
\end{question}

\begin{theorem}[Division algorithm]
Given $a\in\Z$ and $d\in\Z^{+}$, there are \textbf{unique} integers $q$ and $r$ with
\[ a=dq+r,\qquad 0\le r<d. \]
In fact $q=\lfloor a/d\rfloor$ and $r=a\bmod d=a-d\lfloor a/d\rfloor$.
\end{theorem}

\begin{intuition}
Mark the multiples of $d$ on the number line, $\dots,-2d,-d,0,d,2d,\dots$. The number $a$ falls into exactly one ``drawer'' $[dq,\,d(q+1))$. The quotient $q$ says which drawer, and the remainder $r$ says how far into the drawer $a$ sits, somewhere between $0$ and $d-1$. That's why \textbf{the remainder is never negative, even when $a$ is}. It's always measured from the multiple on the \emph{left}.
\end{intuition}

\begin{proof}
\textbf{Existence.} Let $S=\{a-dk : k\in\Z,\ a-dk\ge 0\}$. This set isn't empty. If $a\ge 0$ take $k=0$, and if $a<0$ take $k=a$, which gives $a-da=a(1-d)\ge 0$ since $a<0$ and $1-d\le 0$. By the well-ordering principle $S$ has a smallest element $r=a-dq\ge 0$. If $r$ were at least $d$, then $r-d=a-d(q+1)$ would also be in $S$ and smaller than $r$, which is impossible. Hence $0\le r<d$.

\textbf{Uniqueness.} Suppose $a=dq+r=dq'+r'$ with $0\le r,r'<d$. Then $d(q-q')=r'-r$, so $d\mid(r'-r)$. But $-d<r'-r<d$, and the only multiple of $d$ in that range is $0$. So $r=r'$, and since $d\ne0$ we also get $q=q'$.
\end{proof}

\begin{example}
Take $a=-20$ and $d=3$. The multiples of 3 on either side of $-20$ are $-21$ and $-18$. The one on the left is $-21=3\cdot(-7)$, and $-20$ is one step past it. So $-20=3(-7)+1$, with $q=-7$ and $r=1$.
\end{example}

\begin{mistake}
$-20=3(-6)-2$ is a true equation, but it is \textbf{not} the division, because $r=-2$ is negative. If you ever get a negative remainder, subtract 1 from the quotient and add $d$ to the remainder, turning $(-6,-2)$ into $(-7,1)$.
\end{mistake}

\subsection{Greatest common divisor}

\begin{definition}
For $a,b\in\Z$, not both zero, $\gcd(a,b)=\max\{d\in\Z^+ : d\mid a\ \wedge\ d\mid b\}$. The maximum exists because the set contains $1$ and, by the size property, every element is at most $\max(|a|,|b|)$. We call $a$ and $b$ \textbf{coprime}, or relatively prime, when $\gcd(a,b)=1$.
\end{definition}

\begin{question}
How do you find the gcd without factoring? Factoring large numbers is hard, but dividing them is easy.
\end{question}

\begin{lemma}[Euclid's lemma for the gcd]
If $a=bq+r$, then $\gcd(a,b)=\gcd(b,r)$.
\end{lemma}

\begin{proof}
We prove something stronger, namely that \emph{the two pairs have exactly the same common divisors}. Their largest common divisors then have to agree.
\begin{itemize}
\item If $d\mid a$ and $d\mid b$, then $d$ divides the combination $a-bq=r$, so $d$ is a common divisor of $b$ and $r$.
\item If $d\mid b$ and $d\mid r$, then $d\mid bq+r=a$, so $d$ is a common divisor of $a$ and $b$.
\end{itemize}
The two sets of common divisors are equal, so their maximums are equal too.
\end{proof}

The case $q=1$ is worth remembering on its own, $\gcd(x,y)=\gcd(x,y-x)$. You can also prove it by showing $d_1\le d_2$ and $d_2\le d_1$, which is the same idea looked at through the maximums.

\subsection{Euclid's algorithm}

\begin{intuition}
Draw an $a\times b$ rectangle. Cut off as many $b\times b$ squares as will fit, which is $q_1$ of them, and you're left with a $b\times r_1$ rectangle. Do the same with that one, and keep going. The side of the last square, the one that fills what's left exactly, is the gcd. It's the largest square tile that covers the original rectangle with no gaps, and it's exactly what the Euclid chapter of the module 1 video animates.
\end{intuition}

\begin{theorem}
Apply the division algorithm over and over,
\[
\begin{aligned}
a&=bq_1+r_1, & 0&\le r_1<b\\
b&=r_1q_2+r_2, & 0&\le r_2<r_1\\
&\ \,\vdots\\
r_{k-2}&=r_{k-1}q_k+r_k, & 0&<r_k<r_{k-1}\\
r_{k-1}&=r_kq_{k+1}+0.
\end{aligned}
\]
The process stops, and $\gcd(a,b)=r_k$, the \textbf{last nonzero remainder}.
\end{theorem}

\begin{proof}
The process stops because $b>r_1>r_2>\cdots\ge 0$ is a strictly decreasing sequence of natural numbers, and such a sequence can't go on forever. The lemma then gives $\gcd(a,b)=\gcd(b,r_1)=\cdots=\gcd(r_{k-1},r_k)=\gcd(r_k,0)=r_k$.
\end{proof}

\begin{example}
Here is $\gcd(89,55)$.
\[
\begin{array}{r@{\;=\;}l}
89 & 1\cdot 55+34\\
55 & 1\cdot 34+21\\
34 & 1\cdot 21+13\\
21 & 1\cdot 13+8\\
13 & 1\cdot 8+5\\
8 & 1\cdot 5+3\\
5 & 1\cdot 3+2\\
3 & 1\cdot 2+1\\
2 & 2\cdot 1+0
\end{array}
\qquad\Rightarrow\qquad \gcd(89,55)=1.
\]
Every quotient is 1 because 89 and 55 are consecutive Fibonacci numbers. This is the slowest the algorithm can possibly go, a fact known as Lamé's theorem.
\end{example}

\subsection{The Fundamental Theorem of Arithmetic}

\begin{theorem}[FTA]
Every integer $n>1$ is a product of primes, $n=p_1^{a_1}p_2^{a_2}\cdots p_k^{a_k}$, and this factorization is unique apart from the order of the factors.
\end{theorem}

\begin{proof}[Proof idea]
\emph{Existence} goes by strong induction. If $n$ is prime there's nothing to do. Otherwise $n=ab$ with $1<a,b<n$, and both factors split into primes by the induction hypothesis. \emph{Uniqueness} goes like this. If $p_1\cdots p_r=q_1\cdots q_s$, then $p_1$ divides the product of the $q_j$. By \textbf{Euclid's lemma} from Module 2, which says $p\mid ab$ forces $p\mid a$ or $p\mid b$, the prime $p_1$ divides some $q_j$, and since $q_j$ is prime, $p_1=q_j$. Cancel it and repeat.
\end{proof}

\begin{theorem}[Euclid, there are infinitely many primes]
\end{theorem}
\begin{proof}
Suppose there were only finitely many primes $p_1,\dots,p_k$, and let $N=p_1p_2\cdots p_k+1$, which is bigger than 1. By the FTA, $N$ has a prime factor $p$, and $p$ must be one of the $p_i$. Then $p$ divides both $N$ and $p_1\cdots p_k$, so it divides $N-p_1\cdots p_k=1$. That's impossible for a number bigger than 1.
\end{proof}

\begin{counter}
The proof does \textbf{not} say that $N$ is prime. For example $2\cdot3\cdot5\cdot7\cdot11\cdot13+1=30031=59\cdot 509$. What it says is that the prime factors of $N$ are \emph{new}, meaning they aren't on the original list.
\end{counter}

\textbf{The gcd and lcm from a factorization.} Write $a=\prod p_i^{\alpha_i}$ and $b=\prod p_i^{\beta_i}$, allowing some exponents to be $0$. Then
\[
\gcd(a,b)=\prod p_i^{\min(\alpha_i,\beta_i)},\qquad \lcm(a,b)=\prod p_i^{\max(\alpha_i,\beta_i)},\qquad a\cdot b=\gcd(a,b)\cdot\lcm(a,b).
\]
The last identity holds because $\min(\alpha,\beta)+\max(\alpha,\beta)=\alpha+\beta$ for each prime.

\begin{example}
For $a=2^3\cdot3^2\cdot7=504$ and $b=2\cdot3^4\cdot5=810$ we get $\gcd=2^1\cdot3^2=18$ and $\lcm=2^3\cdot3^4\cdot5\cdot7=22680$. As a check, $504\cdot810=408240=18\cdot22680$.
\end{example}

\begin{sketch}
For Euclid's algorithm, (1) divide, (2) swap the pair $(a,b)$ for $(b,r)$, which has the same common divisors, (3) notice the remainders keep shrinking, so the process ends, and (4) use $\gcd(r_k,0)=r_k$.
\end{sketch}

\begin{thumbrule}
Use Euclid for big numbers. If the numbers are already factored, take the smaller exponents for the gcd and the larger ones for the lcm.
\end{thumbrule}

\subsection*{Module 1 exercises}
\begin{enumerate}[label=\textbf{1.\arabic*}]
\item Find $q$ and $r$ for $a=-47$ and $d=6$.\\
\emph{Solution.} The multiple of 6 just to the left of $-47$ is $-48=6(-8)$, so $-47=6(-8)+1$, with $q=-8$ and $r=1$.
\item Compute $\gcd(252,198)$ with Euclid's algorithm.\\
\emph{Solution.} $252=1\cdot198+54$, $198=3\cdot54+36$, $54=1\cdot36+18$ and $36=2\cdot18+0$, so the gcd is $18$.
\item If $a\mid b$ and $a\mid b+2$, show that $a\mid 2$.\\
\emph{Solution.} $a$ divides the combination $(b+2)\cdot1+b\cdot(-1)=2$.
\item Show that $\gcd(n,n+1)=1$ for every $n$.\\
\emph{Solution.} A common divisor $d$ of $n$ and $n+1$ divides $(n+1)-n=1$, so $d=1$.
\end{enumerate}

\newpage
%======================================================================
\section{Module 2. Bézout, extended Euclid and modular inverses}
%======================================================================

\subsection{Bézout's identity}

\begin{question}
The gcd was defined as the greatest common divisor, which is a \emph{multiplicative} idea. Can we also describe it \emph{additively}, using sums of multiples of $a$ and $b$?
\end{question}

\begin{intuition}
Picture two rulers with no markings, one $a$ cm long and one $b$ cm long. By laying them end to end or measuring the difference between them, the lengths you can reach are exactly the numbers $sa+tb$. With $a=6$ and $b=10$ you only ever get multiples of $2$, and you can get $2$ itself, since $2=2\cdot6-1\cdot10$. The smallest length you can measure is the gcd, and every other length is a multiple of it. In the video this appears as a grid of points $(s,t)$, each labeled with its value $sa+tb$.
\end{intuition}

\begin{theorem}[Bézout]
If $a,b\in\Z$ are not both zero, then $\gcd(a,b)=sa+tb$ for some $s,t\in\Z$. In fact $\gcd(a,b)$ is the \textbf{smallest} positive element of $L=\{sa+tb: s,t\in\Z\}$.
\end{theorem}

\begin{proof}
$L$ has positive elements, for example $a\cdot a+b\cdot b$. By well-ordering it has a smallest one, which we call $g=s_0a+t_0b$.

\emph{$g$ divides $a$.} Divide to get $a=gq+r$ with $0\le r<g$. Then $r=a-gq=a-(s_0a+t_0b)q=(1-s_0q)a+(-t_0q)b$, so $r$ is in $L$ too. If $r$ were positive it would be an element of $L$ smaller than $g$, so $r=0$ and $g\mid a$. The same argument shows $g\mid b$.

\emph{$g$ is the largest.} If $d\mid a$ and $d\mid b$, then $d\mid s_0a+t_0b=g$, so $d\le g$.

So $g$ is a common divisor that is at least as big as any other, which makes it $\gcd(a,b)$.
\end{proof}

\begin{corollary}[The coprimality test]
$\gcd(a,b)=1$ if and only if $sa+tb=1$ for some $s,t\in\Z$.
\end{corollary}
\begin{proof}
One direction is Bézout. For the other, any common divisor $d$ of $a$ and $b$ divides $sa+tb=1$, so $d=1$.
\end{proof}

\begin{mistake}
For a general $g$, getting $sa+tb=g$ does \textbf{not} make $g$ the gcd. For instance $1\cdot 6+1\cdot 10=16$ while $\gcd(6,10)=2$. All you can conclude is $\gcd(a,b)\mid g$, and the converse only holds when $g=1$.
\end{mistake}

\subsection{Three lemmas that follow from Bézout}

\begin{lemma}[Coprime divisibility]\label{lem:coprime}
If $a\mid bc$ and $\gcd(a,b)=1$, then $a\mid c$.
\end{lemma}
\begin{proof}[Proof, the template]
\begin{enumerate}
\item Since $\gcd(a,b)=1$, Bézout gives $s,t$ with $sa+tb=1$.
\item Multiply through by $c$ to get $sac+tbc=c$.
\item $a$ divides $sac=a(sc)$, and it divides $tbc$ because $bc=ak$ by hypothesis, so $tbc=a(tk)$.
\item So $a$ divides the sum $sac+tbc$, which is $c$.\qedhere
\end{enumerate}
\end{proof}
A typical exercise follows exactly the same steps with $d$ in place of $a$, going from $dx+by=1$ to $adx+aby=a$ and then to $d(ax+ky)=a$.

\begin{counter}
Without coprimality the lemma fails. $6\mid 4\cdot 3=12$, but $6\nmid 3$, because the factor $2$ of $6$ is hiding in the $4$.
\end{counter}

\begin{lemma}[Euclid's lemma]
If $p$ is prime and $p\mid ab$, then $p\mid a$ or $p\mid b$.
\end{lemma}
\begin{proof}
If $p\mid a$ we're done. Otherwise $\gcd(p,a)=1$, because the only positive divisors of $p$ are $1$ and $p$ and the second one doesn't divide $a$. The previous lemma then gives $p\mid b$.
\end{proof}

\begin{counter}
It fails for composite numbers. $6\mid 2\cdot 3$, but $6\nmid2$ and $6\nmid3$. This lemma is what makes primes special. It's the key step in the uniqueness part of the FTA, and it comes back with the self-inverse elements in Module 4.
\end{counter}

\begin{lemma}[Coprime divisors]
If $a\mid c$, $b\mid c$ and $\gcd(a,b)=1$, then $ab\mid c$.
\end{lemma}
\begin{proof}
Write $c=ak$. Since $b\mid ak$ and $\gcd(b,a)=1$, Lemma~\ref{lem:coprime} gives $b\mid k$, say $k=b\ell$. Then $c=ab\ell$.
\end{proof}
\begin{counter}
$4\mid 12$ and $6\mid 12$, but $24\nmid 12$. The two divisors share a factor of $2$, and the product counts it twice.
\end{counter}

\begin{lemma}[The gcd of a product]
If $\gcd(a,b)=1$ and $\gcd(a,c)=1$, then $\gcd(a,bc)=1$.
\end{lemma}
\begin{proof}
Bézout gives $ax_1+by_1=1$ and $ax_2+cy_2=1$. Multiplying the two,
\[
1=(ax_1+by_1)(ax_2+cy_2)=a\underbrace{(ax_1x_2+x_1cy_2+by_1x_2)}_{X}+(bc)\underbrace{(y_1y_2)}_{Y}.
\]
$X$ and $Y$ are integers, so the coprimality test gives $\gcd(a,bc)=1$.
\end{proof}
Multiplying two Bézout identities is a trick worth keeping in mind whenever you need coprimality with a product.

\subsection{The extended Euclidean algorithm}

\begin{question}
Bézout tells us $s$ and $t$ exist. How do we actually find them?
\end{question}

\textbf{Back-substitution.} Solve each division for its remainder and substitute from the bottom up, always keeping the expression in terms of the \emph{two most recent remainders}.

\begin{example}[Writing $1$ as a combination of $89$ and $55$]
Starting from the divisions in the previous section,
\[
\begin{aligned}
1 &= 3-1\cdot 2\\
  &= 3-(5-1\cdot 3)= 2\cdot 3-5\\
  &= 2(8-5)-5 = 2\cdot 8-3\cdot 5\\
  &= 2\cdot 8-3(13-8) = 5\cdot 8-3\cdot 13\\
  &= 5(21-13)-3\cdot13 = 5\cdot 21-8\cdot 13\\
  &= 5\cdot21-8(34-21) = 13\cdot 21-8\cdot 34\\
  &= 13(55-34)-8\cdot34 = 13\cdot 55-21\cdot 34\\
  &= 13\cdot55-21(89-55) = 34\cdot 55-21\cdot 89.
\end{aligned}
\]
To check, $34\cdot55=1870$ and $21\cdot 89=1869$, which differ by 1. Another valid answer is $1=34\cdot89-55\cdot55$, since $3026-3025=1$, so Bézout coefficients are \textbf{not unique}. Whenever $(s,t)$ works, so does $(s+kb/g,\ t-ka/g)$. Both answers give the same inverse, because $34$ and $-55$ are congruent mod $89$.
\end{example}

\textbf{The table method.} Start with $s_0=1,t_0=0$ for the row of $a$ and $s_1=0,t_1=1$ for the row of $b$. Each new row comes from
\[
r_j=r_{j-2}-q_{j-1}r_{j-1},\qquad s_j=s_{j-2}-q_{j-1}s_{j-1},\qquad t_j=t_{j-2}-q_{j-1}t_{j-1}.
\]
\emph{Why it works.} Every row satisfies $r_j=s_ja+t_jb$. The first two rows do, since $a=1a+0b$ and $b=0a+1b$, and because everything is linear, the combination ``row $j-2$ minus $q$ times row $j-1$'' keeps the property.

\begin{center}
\begin{tabular}{cccccc}
\toprule
$j$ & $r_j$ & $q_j$ & $s_j$ & $t_j$ & check $s_j\cdot89+t_j\cdot55$\\
\midrule
0 & 89 & -- & 1 & 0 & 89\\
1 & 55 & 1 & 0 & 1 & 55\\
2 & 34 & 1 & 1 & $-1$ & 34\\
3 & 21 & 1 & $-1$ & 2 & 21\\
4 & 13 & 1 & 2 & $-3$ & 13\\
5 & 8 & 1 & $-3$ & 5 & 8\\
6 & 5 & 1 & 5 & $-8$ & 5\\
7 & 3 & 1 & $-8$ & 13 & 3\\
8 & 2 & 2 & 13 & $-21$ & 2\\
9 & \textbf{1} & -- & $\mathbf{-21}$ & $\mathbf{34}$ & 1\\
\bottomrule
\end{tabular}
\end{center}
The $q_j$ column holds the quotient used to build the next row.

\subsection{Modular inverses}

\begin{definition}
$\bar a$ is an \textbf{inverse of $a$ modulo $m$} if $a\bar a\equiv 1\md m$.
\end{definition}

\begin{theorem}
$a$ has an inverse modulo $m$ if and only if $\gcd(a,m)=1$, and in that case the inverse is unique modulo $m$.
\end{theorem}
\begin{proof}
$a\bar a\equiv1\md m$ means $m\mid a\bar a-1$, that is, $a\bar a-1=mk$ for some $k$, or $\bar a\, a+(-k)m=1$. By the coprimality test this happens exactly when $\gcd(a,m)=1$.

For uniqueness, if $a\bar a\equiv a\bar a'\equiv1$, then $\bar a=\bar a(a\bar a')=(\bar a a)\bar a'\equiv\bar a'$.
\end{proof}

\textbf{Computing it.} Extended Euclid gives $sa+tm=1$. Reducing mod $m$ makes the $tm$ term disappear, so $\bar a\equiv s\md m$.

\begin{example}
Reducing $34\cdot55-21\cdot89=1$ mod 89 gives $34\cdot55\equiv1$, so $55^{-1}\equiv34\md{89}$.
\end{example}

\begin{counter}
$2$ has no inverse mod $6$. The row of 2 in the multiplication table mod 6 is $0,2,4,0,2,4$, and 1 never appears, because $2x-1$ is odd while $6x'$ is even.
\end{counter}

\begin{thumbrule}
To decide whether $a^{-1}\bmod m$ exists, compute $\gcd(a,m)$. If it's 1, the coefficient of $a$ in Bézout, reduced mod $m$, is the inverse.
\end{thumbrule}

\subsection*{Module 2 exercises}
\begin{enumerate}[label=\textbf{2.\arabic*}]
\item Write $18=\gcd(252,198)$ as a linear combination of 252 and 198.\\
\emph{Solution.} $18=54-36=54-(198-3\cdot54)=4\cdot54-198=4(252-198)-198=4\cdot252-5\cdot198$.
\item Find $17^{-1}\bmod 60$.\\
\emph{Solution.} $60=3\cdot17+9$, $17=1\cdot9+8$ and $9=1\cdot8+1$. Working back up, $1=9-8=9-(17-9)=2\cdot9-17=2(60-3\cdot17)-17=2\cdot60-7\cdot17$. So $17^{-1}\equiv-7\equiv53$, and indeed $17\cdot53=901=15\cdot60+1$.
\item Show that if $\gcd(a,b)=1$ then $\gcd(a+b,b)=1$.\\
\emph{Solution.} From $sa+tb=1$ we get $s(a+b)+(t-s)b=1$.
\item Does $15$ have an inverse mod $42$?\\
\emph{Solution.} No, because $\gcd(15,42)=3\neq1$.
\end{enumerate}

\newpage
%======================================================================
\section{Module 3. Congruences, the CRT, Fermat and Wilson}
%======================================================================

\subsection{The congruence relation}

\begin{definition}
$a\equiv b\md m$ means $m\mid(a-b)$, or equivalently $a\bmod m=b\bmod m$.
\end{definition}

\begin{intuition}
This is clock arithmetic. On a clock with $m$ hours, $a$ and $b$ point to the same hour. You can also picture coiling the number line into a spring with $m$ positions per turn. Numbers that end up stacked above one another are congruent, and each column of the spring is a \textbf{class} $[r]_m=\{r+km : k\in\Z\}$.
\end{intuition}

\begin{theorem}
Congruence mod $m$ is an equivalence relation, and it splits $\Z$ into exactly $m$ classes $[0],[1],\dots,[m-1]$.
\end{theorem}
\begin{proof}
It is \emph{reflexive} because $m\mid 0$, \emph{symmetric} because $m\mid a-b$ implies $m\mid b-a$, and \emph{transitive} because $m\mid a-b$ and $m\mid b-c$ give $m\mid(a-b)+(b-c)=a-c$. There are $m$ classes because the division algorithm gives each integer exactly one remainder in $\{0,\dots,m-1\}$.
\end{proof}

\subsection{Algebra of congruences}

\begin{theorem}
If $a\equiv b$ and $c\equiv d\md m$, then $a+c\equiv b+d$, $ac\equiv bd$, $a^k\equiv b^k$, and $P(a)\equiv P(b)$ for any polynomial $P$ with integer coefficients.
\end{theorem}
\begin{proof}
For sums, $(a+c)-(b+d)=(a-b)+(c-d)$. For products we \emph{add and subtract} the same term, $ac-bd=ac-bc+bc-bd=c(a-b)+b(c-d)$, and $m$ divides both parts. Powers follow from products by induction, and polynomials are built from sums and products.
\end{proof}

\textbf{In practice} you can reduce mod $m$ \emph{at any point} in a calculation made of sums, products and powers. For example, $27\equiv 9\cdot3\equiv 4\cdot 3=12\equiv 2\md 5$.

\begin{question}
Adding, multiplying and taking powers always work. What about \emph{dividing}?
\end{question}

\begin{theorem}[Cancellation]
If $ac\equiv bc\md m$, then $a\equiv b\md{\dfrac{m}{\gcd(c,m)}}$. In particular, when $\gcd(c,m)=1$ you can cancel $c$ and keep the same modulus.
\end{theorem}
\begin{proof}
Let $g=\gcd(c,m)$ and write $c=gc'$ and $m=gm'$, so that $\gcd(c',m')=1$. The hypothesis says $m\mid c(a-b)$, that is, $gm'\mid gc'(a-b)$, so $m'\mid c'(a-b)$. Since $m'$ and $c'$ are coprime, the coprime divisibility lemma gives $m'\mid a-b$.
\end{proof}

\begin{mistake}
$4x\equiv4y\md 6$ does \textbf{not} imply $x\equiv y\md6$. Try $x=0$ and $y=3$. Then $0\equiv12\md 6$, but $0\not\equiv3\md6$. The correct conclusion is $x\equiv y\md 3$, because $6/\gcd(4,6)=3$.
\end{mistake}

\subsection{Linear congruences $ax\equiv b\md m$}

\begin{theorem}
Let $d=\gcd(a,m)$. The congruence $ax\equiv b\md m$ has a solution if and only if $d\mid b$. When it does, it has \textbf{exactly $d$} solutions that are different mod $m$, namely
\[
x\equiv x_0+k\cdot\tfrac{m}{d}\md m,\qquad k=0,1,\dots,d-1,
\]
where $x_0$ solves $\tfrac ad x\equiv\tfrac bd\md{\tfrac md}$.
\end{theorem}
\begin{proof}
\emph{Necessity.} If $ax-b=mk$, then $b=ax-mk$, and $d$ divides both $a$ and $m$, so $d\mid b$.

\emph{Sufficiency.} If $d\mid b$, divide everything by $d$ to get $a'x\equiv b'\md{m'}$ with $a'=a/d$, $b'=b/d$, $m'=m/d$ and $\gcd(a',m')=1$. Now $a'$ has an inverse and $x_0\equiv (a')^{-1}b'\md{m'}$, so the solutions mod $m'$ are $x_0+km'$. Mod $m=dm'$, the values $k=0,\dots,d-1$ give different remainders, since they're less than $m$ apart, and $k=d$ brings you back to $x_0$.
\end{proof}

\begin{intuition}
When $d>1$, the values $ax \bmod m$ only ever land on multiples of $d$, and they land on each one $d$ times. If $b$ isn't a multiple of $d$ you never reach it, and if it is, you reach it $d$ times.
\end{intuition}

\begin{example}[$55x\equiv13\md{89}$]
Since $\gcd(55,89)=1$ there is exactly one solution. Multiplying by $55^{-1}\equiv34$ gives $x\equiv13\cdot34=442=4\cdot89+86$, so $x\equiv86\md{89}$. As a check, $55\cdot86=4730=53\cdot89+13$.
\end{example}

\begin{example}[A case with $d>1$, $6x\equiv9\md{15}$]
Here $d=\gcd(6,15)=3$ and $3\mid9$, so there are 3 solutions. Dividing by 3 gives $2x\equiv3\md5$. Since $2^{-1}\equiv3\md 5$, $x_0\equiv9\equiv4\md5$, and the solutions mod 15 are $4$, $9$ and $14$. Indeed $24$, $54$ and $84$ all leave remainder $9$ when divided by $15$.
\end{example}

\subsection{The Chinese Remainder Theorem}

\begin{question}
If I know the remainder of $x$ mod $3$ and mod $5$, do I know $x$ mod $15$?
\end{question}

\begin{intuition}
Place the numbers $0,\dots,14$ on a $3\times5$ grid, each one in row $x\bmod3$ and column $x\bmod 5$. Counting up by one moves you down a row and across a column, diagonally, wrapping around at the edges. Since $3$ and $5$ are coprime, the diagonal \emph{passes through all 15 cells without repeating}, so each pair of remainders matches exactly one number mod 15. On a $2\times4$ grid, where the moduli aren't coprime, the diagonal closes up after only $\lcm(2,4)=4$ steps and leaves cells empty, so some systems have no solution. This grid is the main animation in module 3.
\end{intuition}

\begin{theorem}[CRT]
If $m_1,\dots,m_n$ are pairwise coprime, the system $x\equiv a_i\md{m_i}$ for $i=1,\dots,n$ has a unique solution modulo $M=m_1\cdots m_n$.
\end{theorem}

\begin{proof}
\emph{Existence, by construction.} Let $M_i=M/m_i$. $M_i$ is a product of moduli that are coprime to $m_i$, so $\gcd(M_i,m_i)=1$ by the gcd-of-a-product lemma used a few times. That means $y_i=M_i^{-1}\bmod m_i$ exists. Now set
\[ x=\sum_{i=1}^n a_iM_iy_i. \]
Mod $m_j$, every term with $i\ne j$ vanishes because $m_j\mid M_i$, and the $j$-th term is $a_jM_jy_j\equiv a_j\cdot1$. So $x\equiv a_j\md{m_j}$ for every $j$.

\emph{Uniqueness.} If $x$ and $x'$ are both solutions, every $m_i$ divides $x-x'$. Applying the coprime divisors lemma repeatedly, $M\mid x-x'$.
\end{proof}

\begin{intuition}
Think of each term $C_i=M_iy_i$ as a switch. It is $1$ mod $m_i$ and $0$ mod every other modulus. The solution adds up the remainders, each one multiplied by its own switch.
\end{intuition}

\begin{example}[By substitution]
Solve $x\equiv5\md{55}$ and $x\equiv10\md{89}$. The first equation gives $x=55k+5$. Putting that into the second gives $55k\equiv5\md{89}$, so $k\equiv5\cdot34=170\equiv81\md{89}$. With $k=89q+81$ we get $x=55(89q+81)+5=4895q+4460$, so $x\equiv4460\md{4895}$.
\end{example}

\begin{example}[With the formula]
Solve $2x\equiv1\md5$, $x\equiv-1\md7$, $x\equiv4\md9$ and $3x\equiv2\md{11}$. \textbf{Start by solving each one for $x$.} Since $2^{-1}\equiv3\md 5$ the first becomes $x\equiv3$, the second is $x\equiv6\md7$, and since $3^{-1}\equiv4\md{11}$ the last becomes $x\equiv8$. The system is now $x\equiv3,6,4,8$ mod $5,7,9,11$, with $M=3465$.
\begin{center}
\begin{tabular}{ccccc}
\toprule
$m_i$ & $a_i$ & $M_i$ & $M_i\bmod m_i$ & $y_i=M_i^{-1}$\\
\midrule
5 & 3 & 693 & 3 & 2\\
7 & 6 & 495 & 5 & 3\\
9 & 4 & 385 & 7 & 4\\
11 & 8 & 315 & 7 & 8\\
\bottomrule
\end{tabular}
\end{center}
$x\equiv3\cdot693\cdot2+6\cdot495\cdot3+4\cdot385\cdot4+8\cdot315\cdot8=4158+8910+6160+20160=39388\equiv1273\md{3465}$.
As a check, $1273=5\cdot254+3=7\cdot181+6=9\cdot141+4=11\cdot115+8$.
\end{example}

\begin{example}[Moduli that are \emph{not} coprime]
When the moduli share factors, \textbf{split each equation into prime powers}, which is the CRT used in reverse.
\begin{itemize}
\item $\{x\equiv7\md{10},\ x\equiv4\md{35}\}$ turns into $x\equiv1\md2$, $x\equiv2\md5$, $x\equiv4\md5$ and $x\equiv4\md7$. The two conditions mod 5 contradict each other, so \textbf{there is no solution}.
\item $\{x\equiv7\md{10},\ x\equiv2\md{35}\}$ turns into $x\equiv1\md2$, $x\equiv2\md5$, which shows up twice but agrees, and $x\equiv2\md7$. With the CRT and $M=70$ the switches are $C_1=35\cdot1$, $C_2=14\cdot4=56$ and $C_3=10\cdot5=50$, so $x\equiv35+2\cdot56+2\cdot50=247\equiv37\md{70}$.
\end{itemize}
In general, $x\equiv a\md m$ and $x\equiv b\md n$ can be solved together exactly when $\gcd(m,n)\mid a-b$, and the solution is unique mod $\lcm(m,n)$, not mod $mn$.
\end{example}

\begin{sketch}
For the CRT, (1) $\gcd(M_i,m_i)=1$, so $y_i$ exists, (2) $C_i\equiv1\md{m_i}$ and $C_i\equiv0$ for every other modulus, (3) add up $\sum a_iC_i$, and (4) for uniqueness, if every $m_i$ divides $x-x'$ then so does their product.
\end{sketch}

\subsection{Fermat's Little Theorem}

\begin{question}
How do you compute $2^{1000005}\bmod101$ without doing a million multiplications?
\end{question}

\begin{theorem}[FLT]
If $p$ is prime and $p\nmid a$, then $a^{p-1}\equiv1\md p$. For every $a\in\Z$, $a^p\equiv a\md p$.
\end{theorem}

\begin{intuition}
Multiplying by $a$ \emph{shuffles} the remainders $1,2,\dots,p-1$ without repeating or losing any of them. With $p=7$ and $a=3$ you get $3,6,2,5,1,4$, the same six numbers in a different order. The product of all of them is the same before and after the shuffle, but after the shuffle it contains an extra factor of $a^{p-1}$.
\end{intuition}

\begin{proof}
Look at $a,2a,\dots,(p-1)a$ mod $p$.
\begin{itemize}
\item None of them is $\equiv0$. If $p\mid ka$ with $p\nmid a$, Euclid's lemma gives $p\mid k$, which can't happen for $1\le k\le p-1$.
\item They're all different. If $ia\equiv ja$, we can cancel $a$ since $\gcd(a,p)=1$, which gives $i\equiv j$.
\end{itemize}
So they are $p-1$ different nonzero remainders, which means they are $1,\dots,p-1$ in some order. Multiplying them all together,
\[
a^{p-1}(p-1)!\equiv(p-1)!\md p.
\]
Since $\gcd((p-1)!,p)=1$, the factorial cancels and $a^{p-1}\equiv1$. For the second form, multiply by $a$ when $p\nmid a$, and when $p\mid a$ both sides are $\equiv0$.
\end{proof}

\textbf{Handling huge exponents.} Divide the exponent by $p-1$ to get $E=q(p-1)+r$. Then $a^E=(a^{p-1})^q a^r\equiv a^r\md p$. In other words, \emph{the exponent only matters mod $p-1$}.

\begin{example}
For $2^{1000005}\bmod101$, note that $101$ is prime and doesn't divide 2, so $2^{100}\equiv1$. Since $1000005=100\cdot10000+5$, we get $2^{1000005}\equiv2^5=32\md{101}$.
\end{example}

\begin{example}[An inverse from Fermat]
$7^{22}\equiv1\md{23}$, so $7\cdot7^{21}\equiv1$ and $7^{-1}\equiv7^{21}$. Repeated squaring gives $7^2=49\equiv3$, $7^4\equiv9$, $7^8\equiv81\equiv12$ and $7^{16}\equiv144\equiv6$. Since $21=16+4+1$, $7^{21}\equiv6\cdot9\cdot7=378=16\cdot23+10\equiv10$. As a check, $7\cdot10=70=3\cdot23+1$.
\end{example}

\begin{mistake}
Fermat \textbf{only} works for a prime modulus. For $a^{12}\bmod15$ the exponent $14$ is meaningless. Use Euler with $\varphi(15)=8$, or the CRT with the moduli 3 and 5.
\end{mistake}

\begin{counter}
The converse is false. $2^{340}\equiv1\md{341}$, but $341=11\cdot31$ isn't prime, which makes it a pseudoprime. Having $a^{n-1}\equiv1$ does not prove that $n$ is prime.
\end{counter}

\subsection{Wilson's Theorem}

\begin{theorem}[Wilson]
An integer $n>1$ is prime if and only if $(n-1)!\equiv-1\md n$.
\end{theorem}

\begin{intuition}
In $\{1,\dots,p-1\}$ every number has an inverse, and almost all of them pair up with a \emph{different} number. Each pair $a\cdot a^{-1}=1$ cancels out of the product. The only numbers that are their own inverse are the solutions of $a^2\equiv1$, which are $1$ and $p-1$ (see Module 4). So the whole product is $1\cdot(p-1)\equiv-1$. For $p=7$ this reads $6!=1\cdot(2\cdot4)(3\cdot5)\cdot6\equiv1\cdot1\cdot1\cdot6\equiv-1$.
\end{intuition}

\begin{proof}
First suppose $p$ is prime. For $p=2$ we have $1!\equiv-1\md2$. For odd $p$, each $a\in\{2,\dots,p-2\}$ has an inverse different from itself, because $a^2\equiv1$ means $p\mid(a-1)(a+1)$, and by Euclid's lemma that forces $a\equiv\pm1$. So $\{2,\dots,p-2\}$ splits into pairs with product $1$, and $(p-1)!\equiv1\cdot1\cdots1\cdot(p-1)\equiv-1$.

Now suppose $n=ab$ with $1<a<n$. Then $a\mid(n-1)!$. If $(n-1)!\equiv-1\md n$ held, then since $a\mid n$ we would also have $(n-1)!\equiv-1\md a$, and so $a\mid 1$, which is impossible.
\end{proof}

\begin{example}
With $p=101$, Wilson gives $100!\equiv-1\equiv100$. Since $100!=100\cdot99!$ and $100\equiv-1$, we get $-99!\equiv-1$, so $99!\equiv1$. Since $99!=99\cdot98!$ and $99\equiv-2$, we get $-2\cdot98!\equiv1$, that is, $2\cdot98!\equiv-1\equiv100$, and so $98!\equiv50\md{101}$. As a check, $2\cdot50=100\equiv-1$.
\end{example}

\begin{thumbrule}
For a factorial mod a prime $p$, start from Wilson's $(p-1)!$ and peel off the extra factors by writing them as $-1,-2,\dots$
\end{thumbrule}

\subsection*{Module 3 exercises}
\begin{enumerate}[label=\textbf{3.\arabic*}]
\item Solve $8x\equiv12\md{20}$.\\
\emph{Solution.} $d=4$ divides 12, so there are 4 solutions. Dividing by 4 gives $2x\equiv3\md5$, so $x\equiv4\md5$, and the solutions are $4,9,14,19\md{20}$.
\item Solve $x\equiv2\md3$, $x\equiv3\md5$ and $x\equiv2\md7$.\\
\emph{Solution.} $M=105$, with $C_1=35\cdot(35^{-1}\bmod3)=35\cdot2=70$, $C_2=21\cdot(21^{-1}\bmod5)=21\cdot1=21$ and $C_3=15\cdot(15^{-1}\bmod7)=15\cdot1=15$. So $x\equiv140+63+30=233\equiv23\md{105}$.
\item Compute $3^{100}\bmod7$.\\
\emph{Solution.} $3^6\equiv1$ and $100=6\cdot16+4$, so $3^{100}\equiv3^4=81\equiv4$.
\item Compute $15!\bmod17$.\\
\emph{Solution.} $16!\equiv-1$ and $16\equiv-1$, so $-15!\equiv-1$ and $15!\equiv1$.
\item Can $x\equiv3\md6$ and $x\equiv5\md{8}$ hold at the same time?\\
\emph{Solution.} Yes, because $\gcd(6,8)=2$ divides $5-3=2$. The solutions are the numbers that are $3\bmod6$ and $5\bmod 8$, which is $x\equiv21\md{24}$.
\end{enumerate}

\newpage
%======================================================================
\section{Module 4. Euler's $\varphi$ function, Euler's Theorem and RSA}
%======================================================================

\subsection{The $\varphi$ function}

\begin{definition}
$\varphi(n)=\big|\{x : 1\le x\le n,\ \gcd(x,n)=1\}\big|$, the number of remainders that have an inverse mod $n$.
\end{definition}

\begin{theorem}[Formulas for $\varphi$]
\begin{enumerate}
\item $\varphi(p)=p-1$ when $p$ is prime.
\item $\varphi(p^k)=p^k-p^{k-1}=p^k\left(1-\tfrac1p\right)$.
\item If $\gcd(m,n)=1$, then $\varphi(mn)=\varphi(m)\varphi(n)$.
\item If $n=p_1^{a_1}\cdots p_k^{a_k}$, then $\varphi(n)=n\prod_{i}\left(1-\tfrac1{p_i}\right)$.
\end{enumerate}
\end{theorem}
\begin{proof}
(1) Every number from $1$ to $p-1$ is coprime to $p$.
(2) $\gcd(x,p^k)\ne1$ exactly when $p\mid x$. Between $1$ and $p^k$ there are $p^{k-1}$ multiples of $p$, and those are the ones we subtract.
(3) By the CRT, $x\bmod mn\leftrightarrow(x\bmod m,\,x\bmod n)$ is a bijection, and $\gcd(x,mn)=1$ exactly when $\gcd(x,m)=1$ and $\gcd(x,n)=1$. So the invertible remainders mod $mn$ match the pairs made of an invertible remainder mod $m$ and one mod $n$, and there are $\varphi(m)\varphi(n)$ such pairs. On the CRT grid these are the cells that lie in a good row and a good column at the same time.
(4) Combine (2) and (3).
\end{proof}

\begin{example}
$\varphi(119)=\varphi(7)\varphi(17)=6\cdot16=96$, $\varphi(1000)=\varphi(2^3)\varphi(5^3)=4\cdot100=400$, and $\varphi(360)=360\cdot\frac12\cdot\frac23\cdot\frac45=96$.
\end{example}

\begin{mistake}
$\varphi(mn)=\varphi(m)\varphi(n)$ holds \textbf{only} for coprime $m$ and $n$. Already $\varphi(4)=2$ while $\varphi(2)^2=1$.
\end{mistake}

\subsection{Euler's Theorem}

\begin{theorem}[Euler]
If $\gcd(a,n)=1$, then $a^{\varphi(n)}\equiv1\md n$.
\end{theorem}
\begin{proof}
This is Fermat's proof again, using only the invertible remainders. Let $r_1,\dots,r_{\varphi(n)}$ be the remainders coprime to $n$. The products $ar_1,\dots,ar_{\varphi(n)}$ are still coprime to $n$, by the gcd-of-a-product lemma, and they're all different, since $a$ can be cancelled. So they are the same $r_i$ in a new order. Multiplying everything gives $a^{\varphi(n)}\prod r_i\equiv\prod r_i$, and since $\prod r_i$ is invertible it cancels. Fermat is the special case $n=p$.
\end{proof}

The practical upshot is that exponents can be reduced mod $\varphi(n)$.

\begin{example}[$2^{118}\bmod119$]
$\varphi(119)=96$ and $\gcd(2,119)=1$, so $2^{118}=2^{96}\cdot2^{22}\equiv2^{22}$. Next, $2^7=128\equiv9$, $2^{11}=2^7\cdot2^4\equiv9\cdot16=144\equiv25$, and $2^{22}=(2^{11})^2\equiv625=5\cdot119+30\equiv30$. The remainder is $\mathbf{30}$.
\end{example}

\begin{example}[The last three digits of $7^{2407}$]
The last three digits of a number are its remainder mod 1000. Since $\varphi(1000)=400$ and $2407=6\cdot400+7$, we have $7^{2407}\equiv7^7$. With $7^4=2401\equiv401$ and $7^3=343$, $7^7\equiv401\cdot343=137543\equiv543$. The last three digits are $\mathbf{543}$.
\end{example}

\subsection{Self-inverse elements, $a^2\equiv1$}

\begin{question}
How many remainders are their own inverse? Among the real numbers only $\pm1$ are. Is it the same mod $n$?
\end{question}

\begin{theorem}
\begin{enumerate}
\item If $p$ is an odd prime, $a^2\equiv1\md{p^k}$ has exactly 2 solutions, $\pm1$. We proved this for $p$ above, and the same argument works for $p^k$, because $p$ can't divide both $a-1$ and $a+1$ when their difference is 2.
\item $a^2\equiv1\md{2^n}$ has 1 solution when $n=1$, 2 when $n=2$, and \textbf{4} when $n\ge3$, namely $\pm1$ and $2^{n-1}\pm1$.
\item For a general $n$, the CRT says the number of solutions is the product of the counts for each prime power.
\end{enumerate}
\end{theorem}

\begin{proof}
(1) $p\mid(a-1)(a+1)$, so by Euclid's lemma $p\mid a-1$ or $p\mid a+1$. That gives $a\equiv\pm1$, and these are different because $p\ne2$.

(2) Take $n\ge3$. Then $a$ must be odd, so $a-1$ and $a+1$ are consecutive even numbers. Using $\gcd(x,y)=\gcd(x,y-x)$, $\gcd(a-1,a+1)=\gcd(a-1,2)=2$. So at most one of the two factors has more than one factor of 2, and since their product needs at least $n$ of them, one factor is divisible by $2^{n-1}$.
\begin{itemize}
\item If $2^{n-1}\mid a-1$, then $a\equiv1$ or $a\equiv2^{n-1}+1\md{2^n}$.
\item If $2^{n-1}\mid a+1$, then $a\equiv2^{n-1}-1$ or $a\equiv-1\md{2^n}$.
\end{itemize}
For $n\ge3$ these four values are different, and all four work, since $(2^{n-1}\pm1)^2=2^{2n-2}\pm2^n+1\equiv1$ when $2n-2\ge n$.
\end{proof}

\begin{example}[Mod 100]
$100=4\cdot25$. Mod 4 the solutions are $a\equiv1,3$, and mod 25 they're $a\equiv\pm1$, so there are $2\cdot2=4$ solutions, $\{1,49,51,99\}$. Now multiply \textbf{all 40 remainders coprime to 100}. Every one that isn't its own inverse cancels with its partner, which leaves $1\cdot49\cdot51\cdot99\equiv1\cdot(-51)\cdot51\cdot(-1)=51^2=2601\equiv1\md{100}$. This generalizes Wilson, since the product of the invertible elements is $-1$ only when there are exactly two self-inverse elements.
\end{example}

\begin{example}[Mod 3600]
$3600=2^4\cdot3^2\cdot5^2$. The solutions are $\{1,7,9,15\}$ mod 16, $\{1,8\}$ mod 9 and $\{1,24\}$ mod 25, for a total of $4\cdot2\cdot2=16$. With $C_1=225\cdot1$, $C_2=400\cdot7=2800$ and $C_3=144\cdot4=576$, the formula $a\equiv225r_1+2800r_2+576r_3\md{3600}$ gives
\[\{1,199,449,649,1151,1351,1601,1799,1801,1999,2249,2449,2951,3151,3401,3599\}.\]
An exhaustive search confirms that these are exactly the $a<3600$ with $a^2\equiv1$.
\end{example}

\subsection{RSA}

\begin{question}
How can anyone send me an encrypted message using a public key, in a way that nobody except me can read it?
\end{question}

\textbf{Making the keys.}
\begin{enumerate}
\item Pick primes $p$ and $q$, and compute $n=pq$ and $\varphi(n)=(p-1)(q-1)$.
\item Pick $e$ with $1<e<\varphi(n)$ and $\gcd(e,\varphi(n))=1$.
\item Compute $d\equiv e^{-1}\md{\varphi(n)}$ with extended Euclid.
\item Publish $(n,e)$ as the public key. Keep $d$ as the private key, and keep $p$, $q$ and $\varphi(n)$ secret too.
\end{enumerate}
To \textbf{encrypt} compute $C=M^e\bmod n$, and to \textbf{decrypt} compute $M=C^d\bmod n$.

\begin{theorem}[Why it works]
If $0\le M<n$, then $(M^e)^d\equiv M\md n$.
\end{theorem}
\begin{proof}
$ed=1+k\varphi(n)=1+k(p-1)(q-1)$ for some $k$. We first show $M^{ed}\equiv M\md p$. If $p\mid M$ both sides are $0$. If not, Fermat gives $M^{ed}=M\cdot(M^{p-1})^{k(q-1)}\equiv M$. The same holds mod $q$, and since $\gcd(p,q)=1$, the uniqueness part of the CRT gives $M^{ed}\equiv M\md{pq}$.
\end{proof}

\begin{example}[RSA from start to finish with small numbers]
Take $p=5$ and $q=11$, so $n=55$ and $\varphi(n)=40$. Choose $e=3$, which is coprime to 40. Since $40=13\cdot3+1$, we have $1=40-13\cdot3$, so $d\equiv-13\equiv27\md{40}$. As a check, $3\cdot27=81=2\cdot40+1$.

To encrypt $M=7$, compute $C=7^3=343=6\cdot55+13$, so $C=13$.

To decrypt $13^{27}\bmod55$, use the CRT. Mod 5, $13\equiv3$ and $3^{27}=3^{4\cdot6+3}\equiv3^3=27\equiv2$. Mod 11, $13\equiv2$ and $2^{27}=2^{10\cdot2+7}\equiv2^7=128\equiv7$. The number that is $2\bmod5$ and $7\bmod{11}$ is $x=7$, so we get $M=7$ back.
\end{example}

\begin{mistake}
$d$ is the inverse of $e$ \textbf{mod $\varphi(n)$, not mod $n$}. Exponents live mod $\varphi(n)$ because of Euler, while bases live mod $n$.
\end{mistake}

\begin{thumbrule}
If the modulus is prime, use Fermat and reduce the exponent mod $p-1$. If it isn't prime but $\gcd(a,n)=1$, use Euler and reduce mod $\varphi(n)$. If you can factor the modulus, solve piece by piece and combine with the CRT. For a factorial mod a prime, use Wilson.
\end{thumbrule}

\subsection*{Module 4 exercises}
\begin{enumerate}[label=\textbf{4.\arabic*}]
\item Find the last two digits of $3^{2026}$.\\
\emph{Solution.} $\varphi(100)=40$ and $2026\equiv26\md{40}$. Then $3^{5}=243\equiv43$, $3^{10}\equiv43^2=1849\equiv49$ and $3^{20}\equiv49^2=2401\equiv1$, so $3^{26}=3^{20}\cdot3^5\cdot3\equiv1\cdot43\cdot3=129\equiv29$. The last two digits are $\mathbf{29}$.
\item Find every $a$ with $a^2\equiv1\md{15}$.\\
\emph{Solution.} The solutions are $\pm1$ mod 3 and $\pm1$ mod 5, and the four combinations give $\{1,4,11,14\}$.
\item In RSA with $p=5$, $q=11$ and $e=3$, someone takes $d=3^{-1}\bmod55=37$. Does that work?\\
\emph{Solution.} Not in general, since $13^{37}\bmod55=18\neq7$. The problem is that $3\cdot37=111\equiv31\md{40}$, which isn't $1$.
\item Compute $\varphi(2^5\cdot 7^2)$.\\
\emph{Solution.} $\varphi(32)\varphi(49)=16\cdot42=672$.
\end{enumerate}

\newpage
%======================================================================
\section{Decision map and common mistakes}
%======================================================================

\begin{center}
\renewcommand{\arraystretch}{1.35}
\begin{tabular}{p{6.3cm}p{8.3cm}}
\toprule
\textbf{If the problem asks for\dots} & \textbf{use\dots}\\
\midrule
$\gcd(a,b)$ of large numbers & Euclid, the last nonzero remainder\\
$s,t$ with $sa+tb=\gcd$ & extended Euclid, by back-substitution or the table\\
$a^{-1}\bmod m$ & check $\gcd=1$ and use Bézout, or $a^{m-2}$ when $m$ is prime\\
$ax\equiv b\md m$ & $d=\gcd(a,m)$, and if $d\mid b$ divide by $d$ to find all $d$ solutions\\
several congruences with coprime moduli & the CRT, by the formula or by substitution\\
several congruences with shared factors & split into prime powers and look for clashes\\
$a^E\bmod p$ with $p$ prime & Fermat, reducing $E$ mod $p-1$\\
$a^E\bmod n$ with $\gcd(a,n)=1$ & Euler, reducing $E$ mod $\varphi(n)$\\
$a^E\bmod n$ with $n$ composite & the CRT with the factors of $n$\\
$k!\bmod p$ & Wilson, peeling off factors such as $p-1\equiv-1$\\
the last $k$ digits & work mod $10^k$ with Euler, using $\varphi(10^k)=4\cdot10^{k-1}$\\
a proof that $d\mid c$ using coprimality & the Bézout template, $sa+tb=1$ times $c$\\
\bottomrule
\end{tabular}
\end{center}

\subsection*{Common mistakes, all together}
\begin{enumerate}
\item Writing a negative remainder in the division algorithm, when $0\le r<d$ always.
\item Reading $sa+tb=g$ as $g=\gcd$, which only works when $g=1$.
\item Cancelling without checking the gcd. From $ac\equiv bc\md m$ you only get $a\equiv b\md{m/\gcd(c,m)}$.
\item Using Fermat with a composite modulus.
\item Applying the CRT to moduli that share factors. There may be no solution, and when there is one it's unique mod the $\lcm$.
\item Writing $\varphi(mn)=\varphi(m)\varphi(n)$ when $m$ and $n$ aren't coprime.
\item Computing the RSA key $d$ mod $n$ instead of mod $\varphi(n)$.
\item Giving one solution to $ax\equiv b$ when there are $d$ of them mod $m$.
\item Assuming $p_1\cdots p_k+1$ must be prime.
\end{enumerate}

\end{document}
